\documentclass[12pt]{article}
\usepackage[T1]{fontenc}
\usepackage[utf8]{inputenc}
\usepackage{lmodern}
\usepackage{textcomp}
\usepackage{enumitem}
\usepackage{geometry}
\usepackage{graphicx}
\usepackage{amsmath, amssymb}
\geometry{margin=1in}
\begin{document}
\begin{center}
Economics - Undergraduate Level Assessment 9 \
Total Marks: 40 \
Answer all questions.
\end{center}

\section*{Multiple Choice (10 marks)}
\begin{enumerate}[label=\arabic*.]
\item As a result of automatic stabilizers during economic expansions government expenditures
\begin{enumerate}[label=(\alph*)]
    \item remain constant and taxes rise.
    \item fluctuate and taxes rise.
    \item fall and taxes remain constant.
    \item rise and taxes fall.
    \item and taxes rise.
\end{enumerate}
\item If the government regulates a monopoly to produce at the allocative efficient quantity, which of the following would be true?
\begin{enumerate}[label=(\alph*)]
    \item The monopoly would make an economic profit.
    \item The monopoly would break even.
    \item The monopoly would increase its prices.
    \item The deadweight loss in this market would increase.
    \item The deadweight loss in this market would decrease.
\end{enumerate}
\item Which of the following statements about a price ceiling is accurate?
\begin{enumerate}[label=(\alph*)]
    \item A price ceiling will have no effect on the quantity of the good supplied.
    \item A price ceiling will cause a shift in the supply curve for the good.
    \item A price ceiling always results in a surplus of the good.
    \item A price ceiling will increase the total revenue of buyers.
    \item An effective price ceiling must be at a price above the equilibrium price.
\end{enumerate}
\item The United States produces rice in a competitive market. With free trade the world price is lower than the domestic price. What must be true?
\begin{enumerate}[label=(\alph*)]
    \item The United States begins to import rice to make up for a domestic shortage.
    \item The United States begins to export rice to make up for a domestic shortage.
    \item The United States begins to import rice to eliminate a domestic surplus.
    \item The United States begins to export rice to eliminate a domestic surplus.
\end{enumerate}
\item If the Federal Reserve was concerned about the "crowding-out" effect they could engage in
\begin{enumerate}[label=(\alph*)]
    \item expansionary monetary policy by lowering the discount rate.
    \item expansionary monetary policy by selling Treasury securities.
    \item contractionary monetary policy by raising the discount rate.
    \item contractionary monetary policy by lowering the discount rate.
\end{enumerate}
\end{enumerate}

\section*{True / False (10 marks)}
\textit{Answer True or False}
\begin{enumerate}[label=\arabic*.]
\setcounter{enumi}{5}
\item True or False: An autonomous increase in spending refers to an increase in total spending that occurs independently of any rise in income.
\item True or False: Air is free because its marginal utility is zero, indicating it is not a scarce good.
\item True or False: A decrease in the price of aluminum will cause the demand for bicycles to fall due to higher production costs.
\item True or False: A Durbin Watson statistic close to zero suggests that the first order autocorrelation coefficient is exactly zero.
\item True or False: Expectations of surpluses of goods in the future will cause the aggregate demand curve to shift to the left.
\end{enumerate}

\section*{Long Form Response (20 marks)}
\textit{Answer each question with a clear and complete explanation.}
\begin{enumerate}[label=\arabic*.]
\setcounter{enumi}{10}
\item Analyze John's financial situation given a constant marginal propensity to consume (MPC) of 3/4 and a break-even point of $7,000. How much will he need to borrow if his income is $3,000?
\item Discuss how the unique ownership of the only original signed copy of The Wealth of Nations by Adam Smith exemplifies the principles of supply and demand, particularly focusing on the implications of a downward sloping supply curve.
\end{enumerate}

\vfill
\noindent\textit{End of Paper}
\end{document}